\documentclass[border=4pt]{standalone}
\usepackage{amsmath,amssymb,mathtools,xcolor,varwidth}
\begin{document}
\begin{varwidth}{28em}
\large\bfseries
Fix the centre of force $S$ — an immoveable point that pulls on the body.      Now chop the flow of time into equal tiny slices $\Delta t$, and watch the body hop along the      polygonal path $ABCDEF$: in each slice it flies dead-straight by pure inertia (\textbf{Law I}),      then receives one sudden centripetal thump toward $S$ at each vertex. This is Newton's \textit{impulse      approximation} — the method of first and last ratios (\textbf{Lemma V, Cor.~4}) applied not to figures but to      \emph{motion}. If we can show every little triangle $S$ sweeps out has the same area, we will have proved      that area swept is proportional to elapsed time — Kepler's Second Law, freed from ellipses and made true for      \emph{any} central force whatsoever.
\end{varwidth}
\end{document}
